\documentclass[11pt,a4paper]{article}
\usepackage[margin=25mm]{geometry}
\usepackage{amsmath,amssymb,amsthm,hyperref}
\newtheorem{lemma}{Lemma}
\newcommand{\gen}[1]{\langle#1\rangle}
\title{Exact block and orbit deficits at a single crossing}
\author{Human Compression project: bounded L044 proof component}
\date{}
\begin{document}\maketitle
Let $G=\gen{\rho_0,\ldots,\rho_{r-1}}\leq A_n$ be transitive,
with every maximal omission subgroup $G_j$ intransitive. Suppose
$G_i$ has two orbits $A_1,A_2$, and its action on $A_j$ preserves
$m_j\geq2$ blocks of size $k_j\geq2$. Thus
$n=k_1m_1+k_2m_2$. In the single-crossing application these are
the two orbits joined by the unique eligible $i$-edge. The counting
argument below does not require the string property or IP.

Choose a minimal subset $M$ of the distinguished generators of
$G_i$ whose images generate the complete simultaneous action on
the two block systems. Let $c_j$ be the number of $\gen M$-orbits
on $A_j$, and put $c=c_1+c_2$. Minimality makes the block images
of $M$ independent; it does not assert that their omission actions
are intransitive or that their projected marking has IP.

\begin{lemma}[A forest on actual subgroup orbits]\label{forest}
There is a forest on the $c$ orbits of $\gen M$, with precisely
one edge of each distinguished colour not in $M$. In particular
\[
 |M|\leq m_1+m_2-2,\qquad c_j\leq k_j,\qquad
 r-|M|\leq c-1.
\]
\end{lemma}
\begin{proof}
The first bound is the retained intransitive independent-length
bound for a subgroup of $S_{m_1}\times S_{m_2}$, obtained by
successive orbit projections and the normal-quotient normalization
lemma together with Whiston's bound. This statement does not
require an intersection property for the quotient.

The group $\gen M$ is transitive on the $m_j$ blocks of $A_j$.
Each of its point orbits therefore meets every block, so has at
least $m_j$ points. This proves $c_j\leq k_j$.

For every colour $l$ not in $M$, choose an actual $\rho_l$-edge
joining different $G_l$-orbits. Such an edge exists: otherwise
$\rho_l$ preserves every $G_l$-orbit setwise and adjoining it
cannot make the whole group transitive. Contract the point orbits
of $\gen M$. Both ends of the chosen edge lie in different
contracted vertices, since $\gen M\leq G_l$.

The chosen edges form a forest. For if a cycle existed, removing
its $l$-edge would leave a path between its endpoints made of
other-colour edges and paths inside $\gen M$-orbits. All these
steps are performed by elements of $G_l$, contradicting the
choice of the $l$-edge. The same argument excludes parallel edges
as cycles of length two. A forest on $c$ vertices has at most
$c-1$ edges, proving the last inequality.
\end{proof}

Define the nonnegative integer deficits
\[
 \alpha=m_1+m_2-2-|M|,\quad
 \beta=k_1+k_2-c,\quad
 \gamma=c-1-(r-|M|),
\]
and put $D=(k_1-2)(m_1-2)+(k_2-2)(m_2-2)$.

\begin{lemma}[The exact rank deficit]\label{deficit}
One has
\[
 n-1-2r=D-3+2(\alpha+\beta+\gamma).
\]
Consequently the desired bound $r\leq\lfloor(n-1)/2\rfloor$
holds whenever $D+2(\alpha+\beta+\gamma)\geq3$.
\end{lemma}
\begin{proof}
The definitions give
$r=k_1+k_2+m_1+m_2-3-\alpha-\beta-\gamma$.
Substitute this into $n-1-2r$ and expand $D$.
\end{proof}

\begin{lemma}[Saturation forces even block sizes]\label{parity}
If $\alpha=\beta=0$, both $k_1$ and $k_2$ are even.
\end{lemma}
\begin{proof}
The block image has an independent family of size
$m_1+m_2-2$ and exactly two orbits. The accepted intransitive
equality lemma from L047 therefore identifies that image as the
full product $S_{m_1}\times S_{m_2}$.

The equality $c_j=k_j$ forces every $\gen M$-orbit in $A_j$
to have size exactly $m_j$, meeting each block in a single point.
On each such orbit, projection to the block system is an equivariant
bijection. Thus the action of $\gen M$ is the diagonal action of
its block image on $k_j$ copies of the natural $m_j$-point action.
The union of these orbits is the full set of points, so the block
projection is faithful on $\gen M$.

In the full block image choose an element which is a transposition
on the first block system and is the identity on the second. Its
unique lift to $\gen M$ has sign $(-1)^{k_1}$. Since every element
of $G$ is even, $k_1$ is even. Interchanging the systems proves
the claim for $k_2$.
\end{proof}

\begin{lemma}[Exact residual numerical configurations]
If the target rank bound fails, either
\begin{enumerate}
\item $D=0$ and $\alpha+\beta+\gamma\leq1$; each orbit's
block system has either two blocks or blocks of size two; or
\item $D=2$, $\alpha=\beta=\gamma=0$, and, up to interchange,
$(k_1,m_1)=(4,3)$ while the other block system has
$\min(k_2,m_2)=2$ and even $k_2$.
\end{enumerate}
\end{lemma}
\begin{proof}
Lemma~\ref{deficit} restricts a failure to $D\leq2$.
For $D=0$ it gives the first case. If $D=1$ or $2$, all three
deficits must vanish. A summand equal to one requires
$(k_j,m_j)=(3,3)$, excluded by Lemma~\ref{parity}. Thus neither
$D=1$ nor a sum of two ones is possible. A summand equal to two
requires $(k_j,m_j)=(3,4)$ or $(4,3)$; the first is again excluded
by parity. The other summand is zero, giving precisely case two.
\end{proof}

\begin{lemma}[Excluding the broad saturated $D=2$ family]
Restore the string C-group and unique-crossing hypotheses.
The residual case $(k_1,m_1)=(4,3)$, $(k_2,m_2)=(2,m)$
with $m\geq3$ and $\alpha=\beta=\gamma=0$ is impossible.
\end{lemma}
\begin{proof}
Saturation identifies $\gen M$ faithfully with $S_3\times S_m$
on four copies of its first natural action and two copies of
its second. Each primitive factor's active indices form an
interval next to $i$ by Lemma~\ref{interval} below. These
intervals cannot be on the same side: the two-anchor deficit
of Lemma~\ref{same-side} would give $|M|\leq m-1$, whereas
saturation gives $|M|=m+1$. Thus they lie on opposite sides,
and each interval acts trivially on the other block orbit.
They end and start, respectively, at the two neighbors of $i$.
This argument includes $m=3,4$ and does not require an extra
small-marking census.

There are four distinguished generators outside $M$ and $i$:
indeed $c=6$ and $\gamma=0$ give $r-|M|=5$. Only the two
outer neighbors of the component intervals can fail to commute
with $\gen M$. At least two of these four generators therefore
centralize $\gen M$. They fix both crossing endpoints, because
the only indices that can move them other than $i$ are already
in $M$. Centralizing $\gen M$ then forces fixation of the whole
copy containing each endpoint.

The full centralizer on the remaining copies is $S_3$ permuting
the three unused three-point copies, times the trivial permutation
group on the single unused $m$-point copy. Here natural symmetric
actions of degree at least three have trivial internal centralizer;
the two factors' different fixed-point sets also prevent exchange
of copies between the two original orbits. A transposition of
three-point copies is odd on the original points. Hence the
centralizer inside $A_n$ is $A_3$, which has no nonidentity
involutions. It cannot contain either of the asserted commuting
distinguished generators. This contradiction proves the claim.
\end{proof}

\begin{lemma}[An anchored maximal symmetric string]\label{anchored}
Let $m\geq3$, and let $(s_0,\ldots,s_{m-2})$ be an independent
involutory string generating $S_m$. Suppose a point $v$ is fixed
by every $s_j$ except $s_0$. Then the noncommuting graph is the
path in this order, and the subgroup
$T=\gen{s_0,\ldots,s_{m-3}}$ has an orbit containing $v$
of size at least $m-1$. Moreover
$C_{S_m}(T)\cap(S_m)_v=1$.
\end{lemma}
\begin{proof}
For $m\geq5$ use the accepted maximal independent-string
classification from L047, including its degree-five and
 degree-six endpoints. The abstract Coxeter path is ordinary
in the natural action, except possibly for the outer degree-six
marking. The latter's distinguished involutions are all
fixed-point-free, and so it cannot satisfy the hypothesis on
$v$. In the ordinary path $v$ is its first terminal point.
The group $T$ is $S_{m-1}$ on the first $m-1$ points, fixing
the final point. Its centralizer fixing $v$ is trivial.

For $m=3$, $s_1$ is the transposition on the other two points,
and $s_0$ is a transposition joining $v$ to one of them. The
assertions follow directly. For $m=4$, the independent pair
$s_1,s_2$ lies in $(S_4)_v\cong S_3$; it consists of two
distinct transpositions. Relabel so that $v=1$, $s_1=(23)$,
and $s_2=(34)$. Since $s_0$ commutes with $(34)$ and moves
$1$, it is $(12)$ or $(12)(34)$. These give, respectively,
$T=S_{\{1,2,3\}}$ and the transitive dihedral group of order
eight on four points. In the second case a centralizer element
fixing one point of a transitive action is the identity; the
first case is immediate as well. Both displayed strings have
the asserted noncommuting path.
\end{proof}

\begin{lemma}[Primitive block intervals]\label{interval}
Restore the string and unique-crossing hypotheses. If a primitive
block image of $\gen M$ has degree at least three, its active
indices in $M$ form one interval, beginning or ending next to $i$.
\end{lemma}
\begin{proof}
Partition a disconnected active set at a gap. The two resulting
nontrivial commuting normal factors of its primitive image are
transitive. A nonidentity element of either factor, since it
centralizes the other transitive factor, is fixed-point-free.
Every active distinguished generator would therefore move the
block containing the crossing endpoint. Only $i-1$ and $i+1$
can do that: every other generator fixes that endpoint itself.
There would be exactly two active generators, one in each
factor. Each factor would be a transitive group generated by a
single involution, forcing degree two, a contradiction. Thus
the active set is an interval. Transitivity forces one of its
generators to move the crossing block; the interval avoids $i$
and must contain one of its neighbors.
\end{proof}

\begin{lemma}[The maximal block-image two-copy branch]\label{maxcopy}
Restore the string C-group and unique-crossing hypotheses.
Suppose $(k_1,m_1)=(2,a)$ and $(k_2,m_2)=(2,b)$, where
$a,b\geq3$, and $\alpha=0$. Then
$r\leq a+b-1=(n-2)/2$, so the desired rank bound holds.
\end{lemma}
\begin{proof}
The block image is the full $S_a\times S_b$, with a maximal
independent involutory string marking; its lift may have a
nontrivial kernel. By Lemma~\ref{interval}, each factor's active
indices form an interval ending or starting next to $i$.

They cannot both lie on the same side. Suppose they did, on the
right. Choose a minimal subfamily $L\subset M$ generating the
$S_b$ image. Normal-quotient normalization and Whiston's bound
force $|L|=b-1$, leaving $a-1$ independent generators outside
$L$. The image of $L$ satisfies Lemma~\ref{anchored}, with the
crossing block as $v$. Its connected noncommuting path forces
$L$ to occupy the initial interval of indices starting at $i+1$.
Every generator in $M\setminus L$ therefore centralizes its
far-end omission subgroup, or all of $L$, and fixes the crossing
block. Lemma~\ref{anchored} makes its $S_b$ image trivial.
Thus these $a-1$ generators form an independent family in $S_a$,
so generate $S_a$ by Whiston's equality clause. But each fixes
the first orbit's crossing block, because its index is not
$i+1$. This is impossible. The left-sided case is identical.

The active intervals instead lie on opposite sides and are
disjoint. Each acts trivially on the other block orbit, and
independence and the full total length force their lengths to
be $a-1$ and $b-1$. Both images satisfy Lemma~\ref{anchored},
with their near endpoint next to $i$.

Every distinguished generator $x$ outside $M$ and $i$ fixes
both crossing endpoints. It centralizes either both full
intervals, or one full interval and the far-end omission of the
other. A centralizer element of a transitive permutation group
fixing one point is trivial; together with Lemma~\ref{anchored},
this makes both block images of $x$ trivial. It fixes each
crossing block pointwise, since that block has size two.
Commutation transports this fixation throughout the block orbit
of the corresponding centralized subgroup. Both original orbits
are fixed pointwise, except possibly the single far two-point
block on the side where only the omission subgroup is centralized.
Thus $x$ is the identity or a single transposition, both impossible
for a distinguished nonidentity element of $A_n$.

There is consequently no generator outside $M$ and $i$, so
$r=|M|+1=a+b-1$. No existence assertion is made for this
lower-rank configuration. No projected IP or triviality of the
lift's kernel has been assumed.
\end{proof}

\begin{lemma}[Opposite primitive block intervals]\label{twocopy}
Restore the string C-group and unique-crossing hypotheses.
Suppose each orbit has a system of two-point blocks, with numbers
$a,b\geq3$, and the two induced block actions are primitive.
If their active intervals in $M$ lie on opposite sides of $i$,
then $r\leq\lfloor(n-1)/2\rfloor$.
\end{lemma}
\begin{proof}
Here $D=0$. If $\alpha\geq2$ the exact deficit proves the result,
while $\alpha=0$ is Lemma~\ref{maxcopy}. A failure with
$\alpha=1$ therefore has $\beta=\gamma=0$. In particular
$c=4$, and precisely two distinguished generators lie outside
$M$ and $i$. Both neighbors of $i$ belong to $M$.

A distinguished generator $x$ outside $M$ and $i$ cannot
centralize all of $\gen M$. It fixes both crossing endpoints.
In each block system its image centralizes the transitive block
image of $M$ and fixes the crossing block, so its block image
is trivial. It consequently fixes that two-point block pointwise.
Commutation transports this fixation to all blocks, making $x=1$.

The two active intervals are disjoint, each acting trivially on
the other's block orbit. Their independent image markings have
lengths $u\leq a-1$, $v\leq b-1$. Since
$u+v=|M|=a+b-3$, one attains its maximum, say $u=a-1$.
Its image is $S_a$ by Whiston, with the anchored marking of
Lemma~\ref{anchored}.

The two outside generators must occupy the two outer neighboring
indices of the intervals; any generator further away would
centralize all of $M$. Let $x$ be the neighbor on the maximal
$S_a$ side. On the other orbit it centralizes the full transitive
block image and fixes the crossing block, hence has trivial block
image and fixes every point by transporting fixation of the
two-point crossing block. On its own side it centralizes the
far-end omission subgroup and fixes the crossing block.
Lemma~\ref{anchored} makes its block image trivial and propagates
pointwise fixation to at least $a-1$ blocks. Thus it is the
identity or a single odd transposition, a contradiction.
\end{proof}

\begin{lemma}[Two anchored orbits consume two deficits]\label{same-side}
Let $M$ be an independent involutory string in $S_a\times S_b$,
where $a,b\geq3$, with both projections transitive. Suppose $M$ occupies
one interval and only its first generator $t$ can move specified
points $v,w$ in the two images. Then $|M|\leq a+b-4$.
\end{lemma}
\begin{proof}
Transitivity forces $t$ to move both $v$ and $w$. Every generator
other than $t$ and its neighbor fixes $v,w$ and commutes with
$t$, so also fixes the distinct points $v^t,w^t$. Those
$|M|-2$ generators remain independent in the simultaneous image,
and lie in $S_{a-2}\times S_{b-2}$. Normal-quotient normalization
and Whiston's bound give
\[
 |M|-2\leq (a-3)+(b-3).
\]
This includes $a=3$ or $b=3$, when the corresponding remaining
one-point factor contributes zero. The asserted inequality
follows. A one-generator string cannot be transitive in degree
at least three, so the displayed removal of two generators
causes no missing endpoint case.
\end{proof}

\begin{lemma}[The complete two-point-block branch]\label{alltwocopy}
Restore the string C-group and unique-crossing hypotheses.
Suppose both orbit block systems have size two, have primitive
block images, and have at least three blocks. Then
$r\leq\lfloor(n-1)/2\rfloor$.
\end{lemma}
\begin{proof}
Both active sets in $M$ are intervals next to $i$ by
Lemma~\ref{interval}. If they are on opposite sides, use
Lemma~\ref{twocopy}. If they are on the same side, their union
is a single interval. Only its generator next to $i$ can move
either crossing block. Lemma~\ref{same-side} gives $\alpha\geq2$,
and the exact deficit, with $D=0$, proves the result.
\end{proof}

\begin{lemma}[A binary anchored string]\label{binary-anchor}
Let $m\geq4$. Suppose $(t_0,\ldots,t_{m-2})$ is an independent
involutory string in $C_2\times S_m$ whose second projection is
transitive. Suppose only $t_0$ has nonidentity first projection,
and only $t_0$ can move a specified point $v$ of the second
projection. Put $T=\gen{t_0,\ldots,t_{m-3}}$ in the second
projection. Its orbit containing $v$ has size at least $m-1$,
and $C_{S_m}(T)\cap(S_m)_v=1$.
\end{lemma}
\begin{proof}
The final $m-2$ generators have trivial first projection and
fix $v$. Their second projections remain independent. Whiston
therefore makes them generate the full $S_{m-1}$ on the other
points. For $m-1\geq5$ their maximal independent-string marking
is an abstract Coxeter path by L047. An outer degree-six marking
is impossible here: its endpoint omission is transitive on all
six points, so a permutation centralizing it must fix its unique
fixed point $v$. The first generator, which centralizes this
omission and moves $v$, rules this out. In the ordinary marking,
the omission of its first generator is $S_{m-2}$, fixing $v$
and one more point $u$. Its centralizer is generated by $(vu)$,
so the second projection of $t_0$ is $(vu)$. The complete
projected marking is the ordinary anchored Coxeter path.

For $m=4$, the two tail generators are two distinct
transpositions in $S_3$. The direct calculation in
Lemma~\ref{anchored} gives exactly the same two possibilities
there: the first second-projection generator is $(12)$ or
$(12)(34)$, followed by $(23),(34)$. Their initial omission
subgroups are $S_3$ on three points or transitive $D_8$.

It remains to check $m=5$, where the tail is a maximal
independent involutory string of $S_4$. Its two commuting end
involutions cannot both be double transpositions, since then
they lie in the normal four-group and adjoining one involution
cannot generate $S_4$. If they are disjoint single
transpositions, the middle generator must be a transposition
joining their pairs: a double transposition would only generate
a group preserving the partition into those pairs. This is the
ordinary Coxeter path. The remaining case, up to reversal and
relabeling, is
\[
 ((23)(45),(34),(45))
\]
on the four points other than $v=1$. Indeed a commuting single
and double transposition share one factor, and the middle
involution must be a single transposition connecting their two
pairs; a double transposition again preserves that pair partition.
If the displayed tail is reversed, its first omission is
transitive $D_8$, making it impossible for $t_0$ to commute
with that omission and move $1$.

For the displayed order, $t_0$ has second projection $(12)$.
This order contradicts independence in $C_2\times S_5$.
Write $\epsilon$ for the nonidentity element of $C_2$.
The first three generators are
\[
 (\epsilon,(12)),\quad(1,(23)(45)),\quad(1,(34)).
\]
Successive conjugation of the first by the second, third and
second yields $(\epsilon,(13))$, $(\epsilon,(14))$ and
$(\epsilon,(15))$. These lifted star transpositions generate
the graph of the sign homomorphism $S_5\to C_2$.
The third generator $(1,(34))$ lies outside that graph, so the
first three already generate $C_2\times S_5$. The final
$(1,(45))$ would be redundant. Thus only the ordinary tail
survives at $m=5$. In all surviving cases, the stated orbit and
centralizer conclusions follow from the explicit paths or $D_8$.
\end{proof}

\begin{lemma}[A four-point orbit in the two-point-block branch]\label{fourpoint}
Restore the string C-group and unique-crossing hypotheses.
Suppose the two orbits have systems of two-point blocks, with
primitive block images of degrees $2$ and $b\geq4$. Then
$r\leq\lfloor(n-1)/2\rfloor$.
\end{lemma}
\begin{proof}
Here $n=4+2b$ and $D=0$. Only $\alpha=0,1$ can violate the
bound. The degree-$b$ active set in $M$ is an interval by
Lemma~\ref{interval}; assume it starts at $i+1$. The only
possible active indices on the two-block orbit are $i-1,i+1$,
since any block exchange moves its crossing endpoint's block.

Suppose first $i-1\in M$. Its block image is the nonidentity
$C_2$ element and is trivial on the other orbit. The rest of
$M$ is the degree-$b$ interval. Both neighbors of $i$ are in
$M$, so all outside generators fix both crossing endpoints.
An outside generator centralizing $M$ has trivial block images
and fixes the two crossing blocks pointwise; commutation
transports that fixation to every block, a contradiction.
Consequently outside generators can occur only at the two
outer neighbors of $M$.

If $\alpha=1$, failure forces $\beta=\gamma=0$, so exactly
two such generators occur. Consider the outer neighbor on the
single-generator side. It centralizes the complete transitive
degree-$b$ block image and fixes its crossing endpoint, hence
fixes that whole original orbit pointwise. On the four-point
orbit it fixes the crossing block setwise (a permutation of two
blocks fixing one is trivial), and fixes that block pointwise.
Only the other two-point block remains, so the generator is
the identity or one odd transposition, impossible.

If $\alpha=0$, the other interval has $b-1$ generators. Its
degree-$b$ images are independent: any redundancy there could
be lifted to redundancy in $M$ using the pure $C_2$ generator.
They consequently form the anchored maximal $S_b$ marking
of Lemma~\ref{anchored}. The outer neighbor on the single-generator
side is excluded as above. At the other outer neighbor,
commutation with the pure $C_2$ generator fixes the four-point
orbit pointwise, by transporting its fixed crossing block.
Lemma~\ref{anchored} fixes every point on the other orbit
except possibly its far two-point block. Parity again excludes
this neighbor. Thus there is no outside generator; the rank
is $|M|+1=b+1$, already below the claimed upper bound.

Now suppose $i-1\notin M$. Then $M$ is the single right-hand
interval and only its first generator can exchange the two
blocks in the four-point orbit or move the other crossing block.
Removing its first two generators leaves an independent family
fixing both degree-two blocks and two degree-$b$ blocks, of
length at most $b-3$. Thus $|M|\leq b-1$, excluding
$\alpha=0$. With $\alpha=1$, failure again forces
$\beta=\gamma=0$ and two outside generators. Apart from
$i-1$ and the far outer neighbor, any outside generator fixes
both crossing endpoints and centralizes $M$, which is
impossible as above. Hence the far outer neighbor $x$ occurs.

Apply Lemma~\ref{binary-anchor} to the simultaneous block
marking $M$. The generator $x$ fixes both crossing endpoints.
On the four-point orbit its block action is trivial, since it
fixes one of two blocks; it centralizes the first generator
of $M$, which exchanges those blocks, and therefore fixes the
whole orbit pointwise. On the other orbit it centralizes the
far-end omission subgroup of $M$. Lemma~\ref{binary-anchor}
makes its block action trivial and propagates fixation from
the crossing block to at least $b-1$ blocks. Once more it is
the identity or a single odd transposition, a contradiction.
\end{proof}

For $n\geq12$, two systems of two-point blocks cannot both
have degree two. Lemmas~\ref{alltwocopy} and~\ref{fourpoint}
therefore cover every case where maximal nontrivial block
systems on both original orbits have block size two.

\begin{lemma}[One simultaneous swap with two large block pairs]\label{single-swap}
Restore the string C-group and unique-crossing hypotheses.
Suppose $m_1=m_2=2$, $k_1,k_2\geq5$, and $\alpha=1$.
Then $r\leq\lfloor(n-1)/2\rfloor$.
\end{lemma}
\begin{proof}
Write $K=k_1+k_2$. A failure forces $\beta=\gamma=0$ and
$r=K$. The set $M$ consists of one involution $\tau$, which
exchanges both pairs of large blocks. Since it has exactly
$K$ point orbits, all of size two, it is fixed-point-free.
It must be $\rho_{i-1}$ or $\rho_{i+1}$; by reversal assume
the former. Identify each pair of large blocks through $\tau$.
The full block image of $G_i$ is the diagonal $C_2$, generated
by $\tau$.

Let $x=\rho_{i-2}$ if that index exists, and let $C$ be the
remaining distinguished generators of $G_i$ other than $\tau,x$.
They commute with $\tau$. Inside the permutations preserving
both original orbits and inducing the stated diagonal block
image, the centralizer of $\tau$ is
\[
 \operatorname{diag}(S_{k_1}\times S_{k_2})\times\gen\tau.
\]
The projection of $\gen C$ onto the diagonal factor is faithful:
its kernel is contained in $\gen\tau$, whereas IP gives
$\gen C\cap\gen\tau=1$. Write its image as $H$, and regard
$\gen C$ as the graph of a character $\chi:H\to C_2$.
Its marked family remains independent and has IP.

Except possibly $c=\rho_{i+1}$, every member of $C$ fixes both
crossing endpoints. Such a generator has $\chi=0$ and fixes
both corresponding row anchors. If $x$ did not exist, then
$|C|=K-2$, and at least $K-3$ members other than $c$ would be
independent in $S_{k_1-1}\times S_{k_2-1}$, whose bound is
$K-4$. This is impossible. Thus $x$ exists and $|C|=K-3$;
the same bound shows that $c$ exists as well. Whiston's equality
clause now identifies
\[
 H_0=\gen{C\setminus\{c\}}
       =S_{k_1-1}\times S_{k_2-1},
\]
fixing exactly the two row anchors in the two original row sets.

The generator $c$ cannot move both row anchors. If it did, every
member of $C$ other than $c$ and its right neighbor would fix
both anchors and their distinct $c$-images. This independent
family has at least $K-5$ members in
$S_{k_1-2}\times S_{k_2-2}$, whose bound is $K-6$.
It cannot fix both anchors either, since then all $K-3$ members
of $C$ would lie in the subgroup of independent length $K-4$.
After exchanging the two original orbits, it therefore joins
exactly the first anchor to its other rows. A full point
stabilizer together with a permutation moving the stabilized
point generates the full symmetric group. Consequently
\[
 H=S_{k_1}\times S_{k_2-1},
\]
with the second row anchor still fixed. Since $\chi$ vanishes
on $H_0$, it vanishes on $H$: the point stabilizer
$S_{k_1-1}$ contains a transposition and normally generates
$S_{k_1}$, and the second factor already belongs to $H_0$.
Thus $C$ acts diagonally, with no block exchanges.

The prefix of $C$ at indices at most $i-3$ is nonempty.
Otherwise $x$ would centralize all of $C$. It fixes both
crossing endpoints and hence preserves all four large blocks
setwise. On each block of the second orbit, centralizing the
natural $S_{k_2-1}$ fixes its unique fixed row anchor. Then
$\tau,C,x$ preserve the two-point pair over that anchor,
contradicting the transitivity of $G_i$ on this original orbit.
The suffix of $C$, starting at $i+1$, is also nonempty.

These two interval subgroups commute and have trivial
intersection by IP, so give a direct-product decomposition of
$H$. The symmetric factors $S_{k_1}$ and $S_{k_2-1}$ are
centerless and directly indecomposable; as in the accepted
L047 factor lemma, they must be the two interval factors.
The suffix contains $c$, which moves the first anchor, so it
is $S_{k_1}$; the prefix is $S_{k_2-1}$.
The total rank is their maximal independent length, so both
markings attain their separate maxima.

Now $x$ centralizes the entire diagonal $S_{k_1}$. Its action
preserves each large block and hence is trivial on the first
original orbit, because the centralizer of a natural symmetric
group of degree at least three is trivial. On the second orbit
it centralizes the endpoint omission of the prefix marking.
For $k_2\geq6$, the accepted L047 extremal classification applies to
$S_{k_2-1}$, including degrees five and six. In the ordinary
Coxeter path the omission is the natural $S_{k_2-2}$; its
centralizer on the $k_2$ rows is generated by the transposition
of the omitted path endpoint with the permanently fixed row
anchor. In the outer degree-six case its omission is transitive
on those six rows. Its centralizer is trivial: a nontrivial
centralizer of a transitive group is semiregular, so would
contain a fixed-point-free element of order two or three.
The centralizers of those cycle types in $S_6$ have orders
$2^3\cdot3!=48$ and $3^2\cdot2!=18$, respectively; neither
can contain this subgroup of order $120$. Thus its centralizer
on all seven rows is trivial as well.

When $k_2=5$, the direct classification of maximal $S_4$
strings in Lemma~\ref{binary-anchor} shows that the omitted
subgroup is either the natural $S_3$ or transitive $D_8$ on
four of the five rows. In the latter case every centralizer
fixes the unique fixed row anchor. Then $x$, just like $C$,
fixes both points over that anchor, and $\tau$ only exchanges
them. This contradicts the transitivity of $G_i$ on the second
original orbit. The $S_3$ case has the same two fixed rows
and transposition centralizer as the ordinary case above.

In the ordinary case, on the large block containing the
crossing endpoint, $x$ fixes that endpoint and therefore cannot
perform this transposition. On the other large block it may
perform at most that single transposition. Thus $x$ is the
identity or an odd transposition on all $n$ points. Both are
impossible. The outer case gives the identity immediately.
This contradiction proves the rank bound.
\end{proof}

\paragraph{Scope and attribution.}
The two block systems and orbit-contraction strategy come from
Cameron--Fernandes--Leemans--Mixer, \emph{Highest rank of a
polytope for $A_n$}, arXiv:1605.09173v1, Section~5.1, physical
pages~23--25. This formulation uses actual omission-orbit cuts
for its forest and explicitly avoids assuming that a quotient
marking has IP or that mere irredundance implies intransitive
omissions. The complete L047 common theory and its older common
normal-quotient/Whiston premises are required, including the
intransitive equality lemma.

This bounded note proves the broad numerical reduction and the
stated saturation consequences and uniform exclusions above,
including the complete two-point-block case for $n\geq12$
with primitive block images. It does not yet exclude the
two-large-block parts of the residual families other than the
one-swap branch with both block sizes at least five; nor does it prove the finer block drawings
in Propositions~5.2--5.6, or complete Proposition~5.7. No new finite computation is required by these proofs, and no new
formal theorem, original credit or conjecture claim is made.
\end{document}
